\section{Who Will Come to the King's Wedding?}
\label{sec:opening}

The Gospel of this Sunday is a story about an invitation. ``The kingdom of
heaven is likened to a king who made a marriage for his son. And he sent his
servants to call them that were invited to the marriage: and they would not
come.'' The king sends again, with word that the dinner is ready and the
fatlings killed; the invited neglect it, ``one to his farm and another to his
merchandise'', and the rest seize the servants and kill them. The king's
armies destroy the murderers and burn their city. Then the king turns to the
roads: ``Go ye therefore into the highways; and as many as you shall find,
call to the marriage. And his servants going forth into the ways, gathered
together all that they found, both bad and good: and the marriage was filled
with guests'' (Mt 22:2--3, 9--10). In the longer form of the reading the story
goes on. ``The king went in to see the guests: and he saw there a man who had
not on a wedding garment. And he saith to him: Friend, how camest thou in
hither not having on a wedding garment? But he was silent.'' The man is bound
and cast out, ``For many are called, but few are chosen'' (22:11--14).

{\small\noindent Scripture is quoted throughout in the Douay--Rheims version as
revised by Bishop Challoner, a public-domain study translation made from the
Latin Vulgate. It is not the text proclaimed at Mass: the English heard in the
celebration is that of the \work{Lectionary for Mass for Use in the Dioceses of
the United States of America}, which alone is the approved wording. Psalms are
numbered as the Lectionary numbers them, with the Vulgate number, which the
Douay--Rheims and the Latin books use, in parentheses.\par}
\medskip

The first reading has already spread the table. ``And the Lord of hosts shall
make unto all people in this mountain, a feast of fat things, a feast of wine,
of fat things full of marrow, of wine purified from the lees.'' On the same
mountain the Lord destroys the bond that held the peoples and the web over the
nations; ``He shall cast death down headlong for ever: and the Lord God shall
wipe away tears from every face'' (Is 25:6--8). The psalm sings of the
shepherd who prepares a table before his own in the face of their enemies,
anoints the head and fills the cup, and whose mercy follows them until they
dwell in the house of the Lord. Paul, closing his letter to the Philippians,
writes that he knows ``both to be full and to be hungry: both to abound and to
suffer need'' and that he can ``do all things in him who strengtheneth me''
(Phil 4:12--13). Before the Gospel the assembly prays, in a verse made from
Paul's prayer for the Ephesians, to know the hope of its calling.

The Missal's own texts, which serve this Sunday in all three years of the
Lectionary's cycle, frame the readings with a psalm verse, three prayers and a
pair of Communion antiphons. The Entrance Antiphon asks who could stand if the
Lord marked iniquities, and answers that with him is propitiation. The Collect
asks that his grace go before and follow the assembly and keep it intent on
good works. The Prayer over the Offerings asks that through these acts of
devout worship the faithful pass over to heavenly glory, and the Prayer after
Communion that those fed with the Body and Blood of Christ be made partakers of
the divine nature. At Communion the Missal offers either the psalm verse on
the rich who hungered while those who seek the Lord lack no good thing, or
John's promise that when the Lord appears his own will be like him, because
they will see him as he is.

The \work{Ordo lectionum Missae} marks what it singles out in each reading by
the title it sets over it. Over the first reading stands \latin{Faciet Dominus
convivium, et absterget lacrimam ab omni facie}, drawn from Isaiah 25:6 and 8;
over the Gospel, \latin{Quoscumque inveneritis, vocate ad nuptias}, from
Matthew 22:9, a verse that both forms of the Gospel contain; over the second
reading, \latin{Omnia possum in eo, qui me confortat}, from Philippians 4:13.
The introduction to the \work{Ordo} explains that the relation between the
readings of one Mass is shown by the choice of these titles, and that in
Ordinary Time the Old Testament reading is chosen for its relation to the
Gospel while the Apostle and the Gospel are each read in their own course
(\work{Praenotanda} 105--107). The two titles that belong together set a feast
the Lord makes, at which the tear is wiped from every face, beside a wedding to
which whoever is found is called. The letter to the Philippians is read over
four Sundays of Year~A, and this is the last of them.

The Sunday's question is therefore an old one: who comes to God's feast, and
what keeps a person from it? The Fathers and Doctors who expounded these
passages answer along three lines, and each line, carried through every text
of the Mass, yields an interpretation of the whole. The first asks what the
feast is and to whom it is given, and reads the king's feast as the wedding
the Father makes for his Son. For Gregory the Great that wedding is the Church
joined to the Son through the Incarnation; for Jerome, Christ's wedding with a
Church gathered from Jews and gentiles; for Irenaeus, a marriage the one
Father prepared from the beginning and filled from all nations; for Hilary of
Poitiers, the sacrament of the eternal glory to be received at the
resurrection. Jerome and Thomas Aquinas read Isaiah's banquet as a feast of
Christ for the nations, and Aquinas himself sets it beside the king's dinner.
This interpretation must carry the parable's hardest verse, the burnt city,
which several of these Fathers read of a people. The second asks why a guest
gathered from the roads by the king's own servants is cast out, and what the
garment is that he lacked. Augustine and Gregory answer charity, which neither
baptism nor faith guarantees; Aquinas counts charity among the ways of putting
on Christ; Chrysostom, Irenaeus and Jerome answer in other terms; Hilary of
Poitiers gives an answer that stands against Augustine's and Gregory's. They
agree that the garment is required of those the call has gathered, and must be
kept. The third asks why the invited would not come. It finds the reason in
the farm and the merchandise, which Gregory, Hilary and Aquinas read as
earthly labour, worldly ambition and the love of gain, and of which Chrysostom
says that even necessary things must be set below the spiritual; and it reads
the second reading and the first Communion antiphon, by the editor's own
joining, as the freedom the invited lacked. None of the witnesses behind these
interpretations expounds this Mass as a whole. Each interpretation gives every
appointed element a place and closes with four senses of its own, literal,
allegorical, moral and anagogical.

\subsection{The appointed elements at a glance}

\begin{maptable}
Entrance Antiphon & \latin{Si iniquitates observaveris, Domine}; Ps 130
(129):3--4, an adaptation &
If the Lord kept account of iniquities, no one could stand; but with him there
is propitiation. The antiphon closes by naming him the God of Israel.\\
Collect & \latin{Tua nos, quaesumus, Domine, gratia} &
A petition that the Lord's grace always go before and follow the assembly, and
keep it constantly intent on good works.\\
First Reading & Is 25:6--10a &
On this mountain the Lord of hosts makes a feast of fat things and refined wine
for all peoples, destroys the bond over the nations, casts death down for ever
and wipes away tears; in that day his people say, Lo, this is our God.\\
Responsorial Psalm & Ps 23 (22):1--3a, 3b--4, 5, 6; response from v.~6cd &
The Lord shepherds the singer, who will want nothing: pasture and water, the
paths of justice, no fear in the shadow of death; a table prepared, the head
anointed, the cup; mercy following, and a dwelling in the house of the Lord.\\
Second Reading & Phil 4:12--14, 19--20; vv.~15--18 are not read &
Paul knows how to be brought low and to abound, to be full and to be hungry,
and can do all things in him who strengthens him; the Philippians did well to
share his tribulation; God's supplying of all their want from his riches in
glory in Christ Jesus; a doxology.\\
Alleluia verse & cf.\ Eph 1:17--18, recast in the first person plural &
A prayer that the Father of our Lord Jesus Christ enlighten the eyes of our
heart, that we may know the hope of our calling.\\
Gospel, longer form & Mt 22:1--14 &
The king's wedding for his son: the invited refuse and kill the servants, the
king's armies burn their city, the hall is filled from the roads with bad and
good, a guest without a wedding garment is cast out; many are called, few
chosen.\\
Gospel, shorter form & Mt 22:1--10 &
The same parable as far as the filling of the hall, without the guest and the
closing saying.\\
Prayer over the Offerings & \latin{Suscipe, Domine, fidelium preces} &
The faithful's prayers are offered with the sacrificial gifts, that through
these acts of devout worship they may pass over to heavenly glory.\\
Communion Antiphon, first option & \latin{Divites eguerunt et esurierunt};
cf.\ Ps 34 (33):11 &
The rich have known want and hunger, but those who seek the Lord will lack no
good thing.\\
Communion Antiphon, second option & \latin{Cum apparuerit Dominus}; 1~Jn 3:2,
an adaptation &
When the Lord appears, those who are his will be like him, because they will
see him as he is.\\
Prayer after Communion & \latin{Maiestatem tuam, Domine, suppliciter
deprecamur} &
Addressed to the divine majesty: as he feeds the assembly with the Body and
Blood of Christ, so may he make them partakers of the divine nature.\\
\end{maptable}

\noindent The two forms of the Gospel are alternatives, and so are the two
Communion antiphons: the Missal prints the second under \latin{Vel}, and one of
them is used. Every other element is required. The formulary has no Sequence,
Offertory antiphon, proper Preface, prayer over the people or solemn blessing,
and the Lectionary gives no shorter form of the other readings. A claim below
that rests on Matthew 22:11--14 does not hold where the shorter form of the
Gospel is read, and a claim that rests on one Communion antiphon holds only
where that antiphon is chosen; each such claim says so.
